% ============================================================================
% Fichier  : partie2/P02_C12_suivre_cycles_maturite.tex
% Titre    : Suivre : Recommandations, Cycles et Maturité
% Auteur   : MINKA MI NGUIDJOÏ Thierry Emmanuel
% Version  : 1.0 -- novembre 2026
% ============================================================================
% Base : ACSI C09 (D_KL, maturité CMMI, seuils 0,90/0,75/0,55/0,35, niveau
% indéterminé si U ≥ 0,20) et ACSI C18 (second cycle de LiZbiZ en novembre
% 2026 : A 17/2/1, B 9/2/1, C 7/3/2 D_KL ≈ 1,37, D 18/1/1 D_KL ≈ 0,26,
% E 8/4/8, U 0,30 → 0,40, D_KL ≈ 0,05, nouveau prestataire à Garoua et
% Bafoussam depuis septembre 2026).
% Valeurs de premier cycle retenues pour A et E : A 16/3/1, E 8/6/6
% (cohérentes avec µ, U, H d'ACSI C09).
% Sources : A1 (D_KL entre cycles) ; Kullback et Leibler (1951) ; Cover et
% Thomas (2006) ; lissage de Laplace (ajout d'une unité) ; CMMI pour
% l'échelle à cinq niveaux ; IIA (2024) pour le suivi des recommandations.
% Calculs du 06/11/2026 (script de vérification) : D_KL brutes A 0,016,
% B 0,030, C 1,37, D 0,256, E 0,049 ; lissées (+1) A 0,010, B 0,018,
% C 0,794, D 0,175, E 0,036.
% Décision du 06/11/2026 : aucun renvoi à des travaux non publiés (A2).
% Corrigés : corriges/Corr_C12.tex.
% ============================================================================

\chapter[Suivre]{Suivre : Recommandations, Cycles et Maturité}
\label{chap:suivi}
\chaptermark{Suivre, cycles, maturité}

\epigraph{%
    « Une recommandation non suivie\\
    est une assurance perdue. »%
}{Règle du suivi}

\niveauChapitre{Intermédiaire}{%
    Maîtrise du Chapitre~\ref{chap:conclure} (mesures, quadrant,
    observation, rapport)%
}

\begin{boxobjectifs}
\begin{itemize}
    \item Organiser le suivi des recommandations~: états, preuves,
          indicateurs, audit de suivi.
    \item Traiter une recommandation refusée comme un risque accepté.
    \item Comparer deux cycles d'audit par la divergence de
          Kullback-Leibler, connaître ses propriétés et la lisser lorsqu'une
          réponse est absente d'un cycle.
    \item Lire la trajectoire d'un domaine dans le quadrant, et reconnaître
          une régression de la connaissance.
    \item Convertir un taux de conformité en niveau de maturité, et le
          distinguer de la capacité COBIT.
    \item Conduire le second cycle d'audit de LiZbiZ.
\end{itemize}
\end{boxobjectifs}

\bigskip

% ============================================================================
\section*{Question directrice : six mois après, qu'est-ce qui a changé~?}
% ============================================================================

En novembre~2026, AUDIX CONSEIL revient chez LiZbiZ. Le directeur général
affirme que « tout a été corrigé ». Le tableau de suivi de Mme~NKENG
indique 14 recommandations sur 18 « mises en œuvre ». M.~BEYALA a deux
questions~: ces recommandations sont-elles réellement mises en œuvre, et le
système a-t-il changé~? La première relève du suivi~; la seconde, de la
comparaison des cycles.

\medskip
\textbf{Plan du chapitre.}
\begin{enumerate}
    \item Le suivi des recommandations (section~\ref{sec:c12-suivi}).
    \item Comparer deux cycles~: la divergence de Kullback-Leibler
          (section~\ref{sec:c12-dkl}).
    \item Le lissage (section~\ref{sec:c12-lissage}).
    \item Le second cycle de LiZbiZ (section~\ref{sec:c12-lizbiz}).
    \item La maturité (section~\ref{sec:c12-maturite}).
\end{enumerate}

\separateur

% ============================================================================
\section{Le Suivi des Recommandations}
\label{sec:c12-suivi}
% ============================================================================

\subsection{Le tableau de suivi}

Le tableau de suivi indique, pour chaque recommandation, sa priorité, son
responsable, son échéance, son état et la pièce qui prouve sa mise en
œuvre. Les Normes mondiales de l'IIA font du suivi une responsabilité de la
fonction d'audit, qui doit s'assurer que les mesures ont été prises ou que
la direction a accepté le risque de ne pas agir~\cite{iia2024gias}.

\begin{tableaucours}{L{2.8cm}L{5.0cm}L{4.4cm}}{%
    \tableauHeader{État} &
    \tableauHeader{Définition} &
    \tableauHeader{Preuve exigée}
}
Non commencée & Aucune action engagée & Aucune \\
En cours & Action engagée, non achevée & Plan daté, premiers livrables \\
Déclarée mise en œuvre & Le responsable affirme l'avoir achevée & Sa
déclaration écrite \\
Vérifiée & L'auditeur a constaté la mise en œuvre et son effet & Test de
l'auditeur, papier de travail \\
Risque accepté & La direction renonce à agir, par écrit & Décision de
l'instance compétente \\
\end{tableaucours}
\vspace{-0.8em}
\begin{center}{\small\itshape Tableau~12.1. Les états d'une
recommandation}\end{center}

\begin{boxalert}
\textbf{Violation typique.} Compter comme « mises en œuvre » les
recommandations que leurs responsables déclarent achevées.

\smallskip
\textbf{Règle~:} une déclaration de mise en œuvre est une déclaration~: elle
vaut \NEK{} tant que l'auditeur n'a pas testé. Le tableau distingue
« déclarée » et « vérifiée ».
\end{boxalert}

\begin{boxscenario}{Opération LIZBIZ. Le tableau de Mme~NKENG}
\textbf{Faits.} Sur 18~recommandations, le tableau en donne 14
« mises en œuvre ». AUDIX CONSEIL teste les 14~: 9 sont vérifiées, 3 sont
partiellement mises en œuvre (par exemple, la revue trimestrielle des
comptes a été faite une fois, puis abandonnée), 2 n'ont pas de pièce.

\medskip
\textbf{Constatation.} Taux de mise en œuvre déclaré~: $14/18 \approx
78~\%$~; taux vérifié~: $9/18 = 50~\%$.

\medskip
\textbf{Analyse} (confiance élevée). L'écart de 28~points est l'équivalent,
pour le suivi, du biais du choix forcé (Chapitre~\ref{chap:preuve})~: une
déclaration non vérifiée surestime l'avancement.
\end{boxscenario}

\subsection{L'audit de suivi}

\begin{boxprocedure}
\textbf{Tester une recommandation déclarée mise en œuvre}
\begin{enumerate}
    \item Relire l'observation d'origine~: la cause, pas seulement le
          symptôme.
    \item Obtenir la pièce de mise en œuvre (procédure, paramétrage,
          compte rendu).
    \item Tester l'efficacité sur la période écoulée depuis la mise en
          œuvre, comme un contrôle nouveau (Chapitre~\ref{chap:preuve}).
    \item Vérifier que la cause a disparu~: la recommandation de D.2.1
          visait la notification des départs, pas seulement la
          désactivation des cinq comptes.
    \item Mettre à jour l'état (vérifiée, partielle, non vérifiée) et la
          cote du contrôle concerné.
\end{enumerate}
\end{boxprocedure}

\subsection{Une recommandation refusée}

La direction peut refuser une recommandation~: c'est son droit. Mais le
refus doit être écrit, motivé, et pris par l'instance qui fixe l'appétence
(Chapitre~\ref{chap:risque}). Il devient une décision d'accepter un risque~;
il est inscrit au registre et présenté au comité d'audit. L'auditeur ne
« ferme » pas une recommandation refusée sans cette trace.

% ============================================================================
\section{Comparer Deux Cycles}
\label{sec:c12-dkl}
% ============================================================================

\subsection{Pourquoi le taux de conformité ne suffit pas}

Un audit ne s'arrête pas au rapport~: le système évolue, et la mission
suivante doit dire s'il s'est amélioré. Le taux de conformité seul peut
rester stable alors que l'état de connaissance change. Il faut comparer des
\emph{distributions} de réponses.

\subsection{La divergence de Kullback-Leibler}

\begin{boxdef}
\termecle{Divergence de Kullback-Leibler}\index{divergence de
Kullback-Leibler}~: pour deux distributions $p$ (cycle $t+1$) et $q$ (cycle
$t$) sur les réponses $\{\OUI, \NON, \NEK\}$~\cite{kullback1951,minka2026nek},
\[
    D_{\mathrm{KL}}(p \,\|\, q) =
    \sum_{m} p(m) \log_2 \frac{p(m)}{q(m)}.
\]
\end{boxdef}

\begin{tableaucours}{L{3.4cm}L{8.8cm}}{%
    \tableauHeader{Propriété} &
    \tableauHeader{Conséquence pour l'auditeur}
}
Positive ou nulle & $D_{\mathrm{KL}} \ge 0$, nulle si et seulement si les
deux distributions sont identiques (inégalité de
Gibbs~\cite{cover2006elements}) \\
Non symétrique & $D_{\mathrm{KL}}(p\|q) \ne D_{\mathrm{KL}}(q\|p)$ en
général~: le cours mesure toujours le nouveau cycle par rapport à
l'ancien \\
Indéfinie si $q(m) = 0 < p(m)$ & Une réponse absente au premier cycle et
présente au second rend la divergence infinie~: on lisse
(section~\ref{sec:c12-lissage}) \\
En bits & Avec $\log_2$, elle s'exprime en bits, comme l'entropie \\
\end{tableaucours}
\vspace{-0.8em}
\begin{center}{\small\itshape Tableau~12.2. Les propriétés de la
divergence}\end{center}

\begin{boxscenario}{Opération LIZBIZ. Le domaine D six mois plus tard}
\textbf{Faits.} En mai~2026, le domaine~D donnait 13~\OUI{}, 2~\NON{},
5~\NEK{}. En novembre~2026, après les corrections, les mêmes questions donnent
18~\OUI{}, 1~\NON{}, 1~\NEK{}.

\medskip
\textbf{Calcul.} $H$ passe d'environ 1,24 à 0,57~bit, ce qui rejoint la valeur
d'environ 0,6~bit que rapporte l'article A0. La divergence entre les deux
cycles vaut
\[
    0{,}90\log_2\frac{0{,}90}{0{,}65} + 0{,}05\log_2\frac{0{,}05}{0{,}10} +
    0{,}05\log_2\frac{0{,}05}{0{,}25} \approx 0{,}423 - 0{,}050 - 0{,}116
    \approx 0{,}26~\text{bit}.
\]

\textbf{Analyse} (confiance moyenne). Le domaine a changé de région~: de la
conformité opaque à la conformité claire. La divergence le confirme.
\end{boxscenario}

\begin{boiteRemarque}{Et si la conformité ne bouge pas~?}
Supposons une question dont la conformité reste à 0,60 sur deux cycles,
mais dont les \NEK{} passent de 0,30 à 0,05 au profit des \NON{} (de 0,10
à 0,35). Le taux de conformité ne voit rien~; la divergence, elle, vaut
environ 0,50~bit. L'organisation sait désormais ce qu'elle ignorait~:
c'est un progrès réel, même s'il révèle des non-conformités.
\end{boiteRemarque}

\begin{boxoutils}{Calculer la divergence}
\begin{lstlisting}[language=Python]
from math import log2
def dkl(nouveau, ancien, lissage=0):
    a = [x + lissage for x in nouveau]
    b = [x + lissage for x in ancien]
    p = [x / sum(a) for x in a]
    q = [x / sum(b) for x in b]
    return sum(pi * log2(pi / qi) for pi, qi in zip(p, q) if pi > 0)

print(round(dkl([18, 1, 1], [13, 2, 5]), 3))     # 0.256
print(round(dkl([18, 1, 1], [13, 2, 5], 1), 3))  # 0.175 (lissée)
\end{lstlisting}
\end{boxoutils}

% ============================================================================
\section{Le Lissage}
\label{sec:c12-lissage}
% ============================================================================

\subsection{Le problème du zéro}

Au premier cycle, une question a reçu 8~\OUI{}, 0~\NON{}, 4~\NEK{}~; au
second, 10~\OUI{}, 2~\NON{}, 0~\NEK{}. Le terme des \NON{} vaut
$p(\NON)\log_2(p(\NON)/0)$~: la divergence est infinie. Ce n'est pas que le
changement soit infini~; c'est que l'échantillon du premier cycle, petit,
n'a observé aucun \NON{} alors qu'il était possible.

\subsection{Le lissage de Laplace}

On ajoute une unité fictive à chaque réponse des deux cycles avant de
calculer les proportions~: $(8+1, 0+1, 4+1)$ et $(10+1, 2+1, 0+1)$. La
divergence devient finie (environ 0,38~bit). Le lissage réduit un peu toutes
les divergences (D passe de 0,256 à 0,175~bit), d'autant plus que les
effectifs sont petits~; c'est le prix de la comparabilité.

\begin{boxalert}
\textbf{Violation typique.} Lisser certaines comparaisons et pas d'autres,
puis classer les domaines par divergence.

\smallskip
\textbf{Règle~:} on applique le même lissage à toutes les comparaisons
d'une mission, et on l'écrit dans le plan d'audit (onglet PARAMÈTRES du
classeur O4). Sinon, les divergences ne sont pas comparables.
\end{boxalert}

% ============================================================================
\section{Le Second Cycle de LiZbiZ}
\label{sec:c12-lizbiz}
% ============================================================================

\begin{boxlecture}{Cinq domaines, deux cycles}
\begin{tableaucours}{L{2.8cm}cccccc}{%
    \tableauHeader{Domaine} &
    \tableauHeader{Mai} & \tableauHeader{Nov.} &
    \tableauHeader{$U$ mai} & \tableauHeader{$U$ nov.} &
    \tableauHeader{$D_{\mathrm{KL}}$} & \tableauHeader{Lissée}
}
A. Gouvernance & 16/3/1 & 17/2/1 & 0,05 & 0,05 & 0,016 & 0,010 \\
B. Sauvegardes & 8/3/1 & 9/2/1 & 0,08 & 0,08 & 0,030 & 0,018 \\
C. Stocks & 1/10/1 & 7/3/2 & 0,08 & 0,17 & 1,370 & 0,794 \\
D. Accès & 13/2/5 & 18/1/1 & 0,25 & 0,05 & 0,256 & 0,175 \\
E. Exploitation & 8/6/6 & 8/4/8 & 0,30 & 0,40 & 0,049 & 0,036 \\
\end{tableaucours}

\medskip
\textbf{Constatations.}
\begin{itemize}
    \item \emph{C} a le plus changé~: de la non-conformité claire
          ($\mu = 0{,}08$) à la conformité claire ($\mu \approx 0{,}58$,
          $U \approx 0{,}17$). Les correctifs des stocks ont produit un effet
          mesurable.
    \item \emph{D} passe de la conformité opaque à la conformité claire.
    \item \emph{E} a peu changé en divergence (0,05), mais son
          incertitude a \emph{augmenté}~: $U$ passe de 0,30 à 0,40. Un
          nouveau prestataire assure l'exploitation à Garoua et à Bafoussam
          depuis septembre~2026, et les responsables de site ne savent plus
          ce qu'il fait.
    \item \emph{A} et \emph{B} sont stables.
\end{itemize}

\medskip
\textbf{Analyse} (confiance moyenne). La divergence faible de E ne signifie
pas « rien n'a changé »~: elle masque une \textbf{régression de la
connaissance}, que seul $U$ montre. On ne lit pas plus $D_{\mathrm{KL}}$
seule qu'on ne lit $H$ seule.
\end{boxlecture}

\noindent\begin{minipage}{\linewidth}
\begin{center}
\begin{tikzpicture}[scale=1.0]
\fill[green!18] (3,0) rectangle (6,2);
\fill[red!15] (0,0) rectangle (3,2);
\fill[orange!25] (3,2) rectangle (6,5);
\fill[orange!15] (0,2) rectangle (3,5);
\draw[thick, color=CouleurPrimaire] (0,0) rectangle (6,5);
\draw[dashed] (3,0) -- (3,5); \draw[dashed] (0,2) -- (6,2);
\draw[-{Stealth}, very thick, CouleurAccent] (0.5,0.8) -- (3.5,1.7) node[above, font=\scriptsize\bfseries] {C};
\draw[-{Stealth}, very thick, CouleurAccent] (3.9,2.5) -- (5.4,0.5) node[right, font=\scriptsize\bfseries] {D};
\draw[-{Stealth}, very thick, red!70!black] (2.4,3.0) -- (2.4,4.0) node[above, font=\scriptsize\bfseries] {E};
\node[font=\small] at (3,-0.5) {Taux de conformité $\mu$ (0 à 1)};
\node[font=\small, rotate=90] at (-0.5,2.5) {$U \times 10$};
\end{tikzpicture}
\end{center}
\vspace{-0.4em}
\begin{center}{\small\itshape Figure~12.1. Trajectoires de trois domaines
entre mai et novembre~2026}\end{center}
\end{minipage}

\begin{boxmetier}{Auditeur SI}
\textbf{Préparer un second cycle}
\begin{enumerate}
    \item Reprendre la même décomposition et les mêmes questions~: sans
          cela, les distributions ne sont pas comparables.
    \item Interroger autant que possible les mêmes répondants, et noter
          ceux qui ont changé.
    \item Fixer le lissage avant le terrain.
    \item Lire chaque domaine par $\mu$, $U$, $H$ et $D_{\mathrm{KL}}$, et
          tracer sa trajectoire dans le quadrant.
    \item Signaler toute hausse de $U$~: c'est une perte de connaissance,
          souvent due à un changement d'organisation ou de prestataire.
\end{enumerate}
\end{boxmetier}

% ============================================================================
\section{Maturité et Capacité}
\label{sec:c12-maturite}
% ============================================================================

Les directions aiment résumer un domaine par un niveau de 1 à 5. Le cours
propose une conversion du taux de conformité en niveau de maturité, sur
l'échelle à cinq niveaux du modèle CMMI. C'est une \textbf{règle
pédagogique déclarée}, et non une norme.

\begin{tableaucours}{L{2.8cm}cL{3.0cm}L{4.6cm}}{%
    \tableauHeader{$\mu$} &
    \tableauHeader{Niveau} &
    \tableauHeader{Désignation} &
    \tableauHeader{Descripteur}
}
0,90 et plus & 5 & Optimisé & Amélioration continue systématique \\
0,75 à 0,90 & 4 & Quantitatif & Processus mesurés par des indicateurs \\
0,55 à 0,75 & 3 & Défini & Processus standardisés et diffusés \\
0,35 à 0,55 & 2 & Géré & Pratiques documentées et répétables \\
Moins de 0,35 & 1 & Initial & Pratiques ponctuelles, dépendant des
personnes \\
\end{tableaucours}
\vspace{-0.8em}
\begin{center}{\small\itshape Tableau~12.3. Conversion du taux de conformité
en niveau de maturité (convention du cours)}\end{center}

Deux règles encadrent cette conversion~:
\begin{itemize}
    \item \textbf{Pas de niveau pour un domaine opaque.} Si $U \geq 0{,}20$, le
          niveau est \emph{indéterminé}~: un niveau calculé sur des réponses
          inconnues serait trompeur. En mai, chez LiZbiZ, les domaines D et E
          n'ont pas de niveau~; A est au niveau~4, B au niveau~3, C au
          niveau~1. En novembre, D passe au niveau~5, C au niveau~3, A au
          niveau~4, B au niveau~4 ($\mu = 0{,}75$)~; E reste indéterminé.
    \item \textbf{Ce n'est pas la capacité COBIT.} Le cours GCI
          (MSN-2S3-MSI03) évalue la \emph{capacité} d'un processus COBIT par
          niveaux d'activités, avec l'échelle F/L/P/N et la règle de
          l'ISO/IEC~33000. Le niveau de maturité de ce cours mesure la
          conformité constatée d'un domaine. Les deux échelles répondent à des
          questions différentes et ne se convertissent pas l'une dans l'autre.
\end{itemize}

\begin{boxalert}
\textbf{Violation typique.} Annoncer au conseil que « le contrôle d'accès
est passé du niveau 3 au niveau 5 » sans dire qu'en mai le niveau était
indéterminé.

\smallskip
\textbf{Règle~:} une progression de niveau n'a de sens que si les deux
niveaux sont déterminés. En mai, D était opaque~: on dit qu'il est passé de
«~indéterminé~» à «~5~».
\end{boxalert}

\begin{boxmetier}{Contrôleur interne}
\textbf{Tenir le suivi entre deux audits.} Mme~NKENG distingue désormais,
dans son tableau, les recommandations « déclarées » et « vérifiées ». Elle
teste elle-même un échantillon des déclarations chaque trimestre, et
présente au comité d'audit les deux taux.
\end{boxmetier}

\begin{boxmetier}{Responsable de la conformité}
\textbf{Suivre les recommandations à portée légale.} Les recommandations
liées à la loi n°~2024/017 (protection des données) ou au règlement COBAC
ont une échéance qui ne dépend pas de la direction. Le responsable de la
conformité les isole dans le tableau et alerte dès qu'une échéance légale
est menacée.
\end{boxmetier}

\begin{boxinfo}
\textbf{Ce que dit la recherche.} L'article A1~\cite{minka2026nek} propose la
divergence de Kullback-Leibler pour comparer deux cycles d'audit, et l'axiome
d'individuation temporelle d'A0~\cite{minka2026auditologie} pose que l'état
du système à un cycle dépend des interventions du cycle précédent. L'étude
statistique de trajectoires sur de nombreux cycles reste ouverte~; ce
chapitre se limite à la comparaison de deux cycles, qui repose sur des
résultats établis~\cite{kullback1951,cover2006elements}.
\end{boxinfo}

\separateur

\begin{boxresume}
\textbf{Ce qu'il faut retenir de ce chapitre~:}
\begin{itemize}
    \item \textbf{Le suivi} distingue les recommandations déclarées et
          vérifiées~; une déclaration vaut \NEK{} tant qu'elle n'est pas
          testée~; un refus est un risque accepté par écrit.
    \item \textbf{La divergence de Kullback-Leibler} compare deux cycles~;
          elle est positive, non symétrique, et infinie si une réponse
          absente apparaît~: on lisse, toujours de la même façon.
    \item \textbf{La trajectoire dans le quadrant} dit plus qu'un chiffre~:
          une divergence faible peut cacher une régression de la
          connaissance, que seul $U$ révèle.
    \item \textbf{Le niveau de maturité} est une convention déclarée,
          indéterminé pour un domaine opaque, et distinct de la capacité
          COBIT.
\end{itemize}
\end{boxresume}

\begin{boxinfo}
\textbf{Points de contrôle de l'élève.} Je sais~: distinguer « déclarée »
et « vérifiée »~; calculer une divergence, avec et sans lissage~; expliquer
pourquoi elle peut être infinie~; tracer une trajectoire dans le quadrant~;
attribuer un niveau de maturité ou le déclarer indéterminé.
\end{boxinfo}

\bigskip

\begin{boxexo}[title={\ding{45}\quad Exercice~12.1. Les états \niveaubase}]
Attribuez un état du Tableau~12.1 à chacune des recommandations suivantes~:
(a) la DRH écrit que la notification des départs est en place~; (b)
l'auditeur a testé la notification sur les 12 départs d'octobre, tous
traités sous 48~heures~; (c) le conseil a décidé par écrit de ne pas
financer le site de secours de Garoua~; (d) aucun document n'existe sur
l'inventaire tournant.
\end{boxexo}

\begin{boxexo}[title={\ding{45}\quad Exercice~12.2. Les deux taux \niveaubase}]
Sur 25~recommandations, 20 sont déclarées mises en œuvre~; l'auditeur en
vérifie 12, en trouve 4 partielles et 4 sans pièce. Calculez les deux taux
et rédigez la phrase du rapport de suivi.
\end{boxexo}

\begin{boxexo}[title={\ding{45}\quad Exercice~12.3. Une divergence \niveaumoyen}]
Sur le domaine « gestion des changements », l'audit interne de la BANQUE
SANAGA obtient 8~\OUI{}, 4~\NON{}, 8~\NEK{} en 2026, puis 12~\OUI{},
6~\NON{}, 2~\NEK{} en 2027.
\begin{enumerate}[(a)]
    \item Calculez $\mu$, $U$ et $H$ pour les deux cycles, et situez le
          domaine dans le quadrant à chaque cycle.
    \item Calculez la divergence de Kullback-Leibler (posez le calcul).
    \item Le domaine s'est-il amélioré~? Argumentez sans vous limiter au
          taux de conformité.
    \item Quel niveau de maturité attribuer en 2026 et en 2027~?
\end{enumerate}
\end{boxexo}

\begin{boxexo}[title={\ding{45}\quad Exercice~12.4. Non symétrique \niveaumoyen}]
Calculez $D_{\mathrm{KL}}$ du domaine~D dans les deux sens (novembre par
rapport à mai, et mai par rapport à novembre). Pourquoi le cours fixe-t-il un
sens~?
\end{boxexo}

\begin{boxexo}[title={\ding{45}\quad Exercice~12.5. Le zéro \niveaumoyen}]
Reprenez l'exemple de la section~\ref{sec:c12-lissage} (8/0/4 puis
10/2/0). Calculez la divergence lissée (ajout d'une unité) en posant le
calcul. Que deviendrait-elle avec un lissage de 0,5~?
\end{boxexo}

\begin{boxexo}[title={\ding{45}\quad Exercice~12.6. La régression \niveaumoyen}]
Expliquez au directeur général de LiZbiZ, en cinq lignes, pourquoi le
domaine~E « n'a pas bougé » selon la divergence mais s'est dégradé selon
l'auditeur. Quelle recommandation formulez-vous~?
\end{boxexo}

\begin{boxexo}[title={\ding{45}\quad Exercice~12.7. Constance sans progrès~? \niveaumoyen}]
Vérifiez par le calcul l'exemple de la remarque « Et si la conformité ne
bouge pas~? » ($0{,}60/0{,}10/0{,}30$ puis $0{,}60/0{,}35/0{,}05$).
\end{boxexo}

\begin{boxexo}[title={\ding{45}\quad Exercice~12.8. La maturité \niveaubase}]
Attribuez un niveau de maturité, ou déclarez-le indéterminé, pour~: (a)
$\mu = 0{,}82$, $U = 0{,}10$~; (b) $\mu = 0{,}82$, $U = 0{,}25$~; (c)
$\mu = 0{,}30$, $U = 0{,}05$. Pourquoi (b) n'a-t-il pas de niveau, alors que
sa conformité est la même que (a)~?
\end{boxexo}

\begin{boxexo}[title={\ding{45}\quad Exercice~12.9. L'audit de suivi \niveauexpert}]
La recommandation « revue trimestrielle des comptes à privilèges » est
déclarée mise en œuvre. Rédigez le papier de travail de l'audit de suivi~:
objectif, période, test, preuve, critère de vérification, cote, et ce qui
change dans la fiche D.4.2.
\end{boxexo}

\begin{boxexo}[title={\ding{45}\quad Exercice~12.10. Le rapport du second cycle \niveauexpert}]
Rédigez la synthèse d'une page du rapport du second cycle de LiZbiZ~:
évolution par domaine (région, niveau de maturité, divergence), taux de
mise en œuvre vérifié, signalement de la régression du domaine~E, nouvelles
priorités.
\end{boxexo}

\bigskip

\begin{boxinfo}
\textbf{Pour aller plus loin.}
\begin{itemize}
    \item S.~Kullback et R.~Leibler, \textit{On Information and
          Sufficiency}~\cite{kullback1951}~: l'article qui définit la
          divergence.
    \item T.~Cover et J.~Thomas, \textit{Elements of Information
          Theory}~\cite{cover2006elements}~: les propriétés de la
          divergence, démontrées.
    \item IIA, Normes mondiales d'audit interne~\cite{iia2024gias}~: le
          suivi des recommandations.
\end{itemize}
\end{boxinfo}

\begin{boxinfo}
\textbf{Transition.} La Partie~\ref{part:methode} a donné la méthode, de la
préparation au suivi. La Partie~\ref{part:normatif} en présente le cadre~:
les normes professionnelles qui l'encadrent, les référentiels qui lui
fournissent ses critères, le droit camerounais et africain, et la
déontologie. Le Chapitre~\ref{chap:normes} ouvre cette partie.
\end{boxinfo}
